 % \documentclass[10pt,respuestas,a4]{examenEPN}
\documentclass[10pt,a4]{examenEPN0}
\examen{Prueba 02}
\usepackage{comandosEPN}

\usepackage{empheq}

\newcommand{\itemcp}{\stepcounter{enumi}\item[ \textcolor{red}{\theenumi.$^*$}] }

\newcommand{\iitemcp}{\stepcounter{enumii}\item[ \textcolor{red}{\theenumii)$^*$}] }


\materia{Álgebra Lineal}
%\autor{Departamento de Formación Básica}
\autor{GR \hspace{1cm}.}
%\fecha{Semestre 2026-A}
\fecha{Nombre: \hspace{9cm} Cédula:}
\logouno[1.6cm]{Logo_EPN.eps}
\logodos[2.6cm]{Logo_DFB.eps}

\newcommand{\la}{\lambda}

%Datos
\newcommand{\datosAlgebra}[2]{
\EnExamen
{
\vspace{5mm}
    \noindent
    Nombres: \hrulefill\
    Carrera: \rule{4.5cm}{0.4pt}\
    \ifthenelse{\equal{#1}{}}%
    {Paralelo: \underline{\parbox{1.8cm}{\parbox{1cm}{}}}\\[3mm]}
    {Paralelo: \underline{\parbox{1.8cm}{\centering #1}}\\[3mm]}
    C\'edula: \rule{5cm}{0.4pt}\
    Firma: \hrulefill\
    Fecha: \underline{\parbox{2.1cm}{\centering #2}}
}}


% \newcommand{\indicacionesPrueba}{%
% \EnExamen{
% \begin{indicaciones}
%     \item La prueba tiene una duraci\'on  de \textbf{120 minutos.}
%     \item No se permite la utilizaci\'on de dispositivos electr\'onicos (celular, tablets, calculadora, etc.)
%     \item Leer detenidamente los enunciados, \textbf{justificar la respuesta para acceder al puntaje total}.
%     \item La prueba es completamente individual. Los estudiantes deberán dejar sus teléfonos celulares guardados y no podrán tenerlos consigo (por ejemplo, en los bolsillos de su ropa) durante toda la evaluación. En caso de detectarse algún acto de deshonestidad académica o el incumplimiento de esta norma, no podrán seguir rindiendo la prueba y su calificación será \textbf{cero}.
% \end{indicaciones}
% \pagestyle{empty}
% }}

\newcommand{\indicacionesPrueba}{%
\EnExamen{
\begin{indicaciones}
    \item La prueba tiene una duraci\'on  de \textbf{60 minutos.}
    \item Leer detenidamente los enunciados, \textbf{justificar la respuesta} para acceder al puntaje total.
    \item No se permite la utilizaci\'on de dispositivos electr\'onicos (celular, tablets, calculadora, etc.)
    \item La prueba es completamente individual
    \item No se permite portar teléfonos celulares, relojes inteligentes ni otros dispositivos electrónicos durante el examen, aun cuando permanezcan en los bolsillos. Estos dispositivos deberán guardarse en la mochila antes del inicio de la evaluación y esta deberá mantenerse cerrada durante todo el examen.
    \item He leído las indicaciones y acepto que el incumplimiento de estas se penaliza con \textbf{cero}. Firma:...................................

    
\end{indicaciones}
\pagestyle{empty}
}}

\newcommand{\indicacionesExamen}{
\EnExamen{\begin{indicaciones}
    \item La evaluación tiene una duraci\'on  de \textbf{120 minutos.}
    \item No se permite la utilizaci\'on de dispositivos electr\'onicos (celular, tablets, calculadora, etc.)
    %\item Leer detenidamente los enunciados, justificar la respuesta para acceder al puntaje total.
    \item El examen es completamente individual. En caso de detectarse alg\'un acto de deshonestidad acad\'emica no podr\'a seguir rindiendo la prueba y su calificaci\'on ser\'a \textbf{cero}.
\end{indicaciones}
}\pagestyle{empty}
}

%% Imprimir las una hoja A4, dos veces en A5
\newcommand{\imprimirDoble}
{%
\usepackage{pgfmorepages}
\pgfmorepagesloadextralayouts
\pgfpagesuselayout{repeated 2-up}[a4paper,landscape]
\thispagestyle{empty}
}


%Este sí se puede comentar
\newenvironment{solucion}{\color{blue}\proof[Solución]}{\endproof\color{black}}




% \renewcommand{\@autor}{
% ---\\
% \phantom{.} 
% ...}
% \makeatother

\begin{document}
\encabezado
\datosAlgebra{GR6M}{06/11/2026}
%Poner el grupo y la fecha, o dejar las llaves vacías.
\indicacionesPrueba
\begin{preguntas}


%%%%%%%%%%%%%%%%%%%%%%%%%%%%%%%%%%%%%%%%%
\item  Considere el sistema lineal homogéneo $Ax = 0$, donde $A$ es una matriz cuadrada de orden $4 \times 4$. Si se sabe que $\det(A) = 0$, ¿qué se puede afirmar con certeza sobre el conjunto solución del sistema?

Justifique su respuesta \puntaje{2}
    \begin{enumerate}[label=(\alph*), itemsep=2pt]
        \item El sistema es inconsistente (no tiene solución).
        \item El sistema tiene únicamente la solución trivial $x = 0$.
        \item El sistema tiene infinitas soluciones.
        \item El sistema tiene exactamente cuatro soluciones distintas.
    \end{enumerate}   
    \newpage
%%%%%%%%%%%%%%%%%%%%%%%%%%%%%%%%%

\item Un centro de datos ejecuta simultáneamente tres tipos de
procesos informáticos:

\begin{itemize}
\item Procesos de análisis de datos (A),
\item Procesos de respaldo y almacenamiento de información (B),
\item Procesos de inteligencia artificial para reconocimiento de patrones (C).
\end{itemize}
Cada proceso consume recursos computacionales del servidor.
Un proceso de análisis de datos utiliza 2 unidades de CPU,
1 unidad de memoria RAM y 1 unidad de almacenamiento.
Un proceso de respaldo utiliza 3 unidades de CPU y 2 unidades
de memoria RAM. Por su parte, un proceso de inteligencia
artificial utiliza 2 unidades de CPU, 1 unidad de memoria RAM
y 2 unidades de almacenamiento. Durante una hora de operación se registró un consumo total de 21 unidades de CPU, 12 unidades de memoria RAM y
10 unidades de almacenamiento.

\begin{enumerate}
\item Defina las tres variables, plantee el sistema de ecuaciones
lineales que modela la situación y escriba su representación
matricial $Ax=b$, identificando $A$, $x$ y $b$. \puntaje{3}

\item Resuelva el sistema mediante la regla de Cramer y determine los valores de las tres incógnitas
en el contexto del problema. \puntaje{3}
\end{enumerate}
\newpage
%%%%%%%%%%%%%%%%%%%%%%%%%%%%%%%%%%%%%%%
\item Dadas las matrices: \puntaje{6}
    \[A = \left(\begin{array}{cc}
         \alpha &1  \\
         \beta& 1 
    \end{array}\right),\quad B = \left(\begin{array}{cc}
         \alpha &-1  \\
         -\beta& 2 
    \end{array}\right)\]
    donde \(\alpha,\beta\in\R\) . Explique por qué B no es la inversa de A.


\newpage
%%%%%%%%%%%%%%%%%%%%%%%%%%%%%%%%%%%%%%%
\item (Ejercicio del deber)    Determine la inversa de $B$ utilizando el método de Gauss-Jordan, donde \puntaje{6}
$$B=\begin{pmatrix}1&0&2\\2&-1&3\\4&1&8\end{pmatrix}.$$

Una vez calculada $B^{-1}$, utilícela para resolver el sistema matricial $B\begin{pmatrix}x\\y\\z\end{pmatrix} = \begin{pmatrix}2\\-1\\4\end{pmatrix}.$


    
\end{preguntas}
\final
\end{document}